\documentclass{beamer}
\usetheme{metropolis}

\title{Domain-Specific Financial LLMs}
\author{Juan F. Imbet}
\institute{Paris Dauphine -- PSL University}
\date{}

\begin{document}

\begin{frame}
  \titlepage
\end{frame}

\section{Why Domain-Specific LLMs?}

\begin{frame}{Why Domain-Specific LLMs?}
  \begin{itemize}
    \item Financial text is dense with domain-specific jargon, numerical tables, and regulatory language
    \item General-purpose models lack grounding in financial concepts and conventions
    \item Errors in financial NLP tasks can have significant economic consequences
    \item Domain-specific pre-training substantially improves performance on financial benchmarks
  \end{itemize}
\end{frame}

\section{A Taxonomy of Financial LLMs}

\begin{frame}{A Taxonomy of Financial LLMs}
  \begin{itemize}
    \item BERT-family: FinBERT (sentiment), SEC-BERT (filings), FinBERT-tone (tone analysis)
    \item Decoder-family: BloombergGPT, FinGPT, FinMA, InvestLM
    \item Models differ in training data, architecture size, and target tasks
    \item Open-weight vs.\ proprietary models present different deployment trade-offs
  \end{itemize}
\end{frame}

\section{Pre-training Strategies for Finance}

\begin{frame}{Pre-training Strategies for Finance}
  \begin{itemize}
    \item Domain-Adaptive Pre-training (DAPT): continue pre-training a general model on financial corpora
    \item Corpus composition critically affects downstream task performance
    \item Key sources: Bloomberg terminal archives, SEC EDGAR filings, financial news wires
    \item Tokenizer adaptation may be needed to handle financial notation and table structures
  \end{itemize}
\end{frame}

\section{Benchmark Comparisons}

\begin{frame}{Benchmark Comparisons}
  \begin{itemize}
    \item FinQA tests numerical reasoning over financial reports
    \item FLUE (Financial Language Understanding Evaluation) covers sentiment, NER, and classification
    \item BBH-Finance probes broader financial reasoning in large models
    \item FinLLMs outperform on narrow extraction tasks; general models remain competitive on reasoning
  \end{itemize}
\end{frame}

\section{Practical Deployment Considerations}

\begin{frame}{Practical Deployment Considerations}
  \begin{itemize}
    \item Licensing constraints vary widely: open-weight, commercial, and restricted models
    \item Data governance requirements (GDPR, MiFID II) affect hosting and inference choices
    \item Hybrid strategies combine a cheap general model with a specialized FinLLM for complex queries
    \item Cost-performance trade-offs must be evaluated per use-case at deployment time
  \end{itemize}
\end{frame}

\end{document}
